\clearpage
\begin{summarybox}
  \textbf{\textcolor{CSummary}{快记｜一句话}}

  \textbf{对数判正负：先验条件，再看 $1$。底真同侧为正，异侧为负；真数为 $1$ 就是零。}

  \[
    \boxed{\ \log_a b>0\iff(a-1)(b-1)>0\ },\qquad
    \boxed{\ \log_a b<0\iff(a-1)(b-1)<0\ }.
  \]
  条件：$a>0,\ a\neq1,\ b>0$。

  \medskip
  \textbf{\textcolor{CSummary}{快懂｜为什么？}}

  $\log_a b=x$ 就是 $a^x=b$。因为 $a^0=1$，所以 $1$ 是分界点。底数大于 $1$ 和小于 $1$ 时，正负指数产生的大小变化恰好相反。
  \begin{center}
  \begin{tikzpicture}[x=1cm,y=.72cm,>=stealth]
    \draw[->,thick] (0,0)--(9,0);
    \draw[very thick,teal!75!black] (.4,0)--(4.3,0);
    \draw[very thick,orange!85!black] (4.3,0)--(8.6,0);
    \fill (4.3,0) circle (2pt) node[below] {$1$};
    \node[above,teal!75!black] at (2.25,.25) {$0<t<1$};
    \node[above,orange!85!black] at (6.55,.25) {$t>1$};
    \node[below] at (4.4,-.54) {比较 $a$、$b$ 在 $1$ 的哪边};
  \end{tikzpicture}

  \end{center}

  \textbf{\textcolor{CSummary}{快用｜考试三步}}
  \begin{enumerate}[leftmargin=*,itemsep=0.2em,topsep=0.2em]
    \item \textbf{先验：}底数合法、真数大于 $0$。
    \item \textbf{再比：}底数和真数分别与 $1$ 比较。
    \item \textbf{最后判：}同侧正、异侧负、真数为 $1$ 则为零；解不等式时别漏定义域。
  \end{enumerate}
  \textbf{速测：}$\log_2\frac14=-2<0$；$\log_{\frac12}\frac14=2>0$；$\log_31=0$。
\end{summarybox}
